\documentclass[journal]{IEEEtran}

\usepackage{amsmath,amssymb,amsfonts,bm}
\usepackage{graphicx}
\usepackage{cite}
\usepackage{url}
\usepackage{booktabs}
\usepackage{multirow}
\usepackage{microtype}
\usepackage{algorithm}
\usepackage{algpseudocode}
\usepackage{subcaption}
\usepackage[pagebackref=false,breaklinks=true,colorlinks=true,linkcolor=blue,citecolor=blue,urlcolor=blue]{hyperref}

% Math command shortcuts
\newcommand{\R}{\mathbb{R}}
\newcommand{\DTwoRO}{\text{D}^2\text{RO}}
\newcommand{\SWDGO}{\text{SW-DGO}}

\graphicspath{{figures/}}

\begin{document}

\title{Socially-Weighted Distributed Graph Optimization ($\text{D}^2\text{RO}$) for Autonomous Multi-Agent Service Fleets in Crowded Environments}

\author{Polla~Fattah,~\IEEEmembership{Senior Member,~IEEE,}
        and~Research~Collaborators%
\thanks{P. Fattah is with the Department of Computer Science and Robotics, Faculty of Engineering and Applied Sciences. (e-mail: polla.fattah@research.org).}}

\markboth{IEEE Transactions on Robotics,~Vol.~42, No.~1, 2026}%
{Fattah \MakeLowercase{\textit{et al.}}: Socially-Weighted Distributed Graph Optimization ($\text{D}^2\text{RO}$)}

\maketitle

\begin{abstract}
The continuous routing of autonomous service fleets—such as retail shopping carts (\textit{Int-Cart}), clinical hospital pushchairs, and airport luggage trolleys—in crowded, human-shared public environments poses fundamental multi-agent coordination challenges. Existing Multi-Agent Path Finding (MAPF) and reactive collision avoidance methods suffer from local potential minima traps in orthogonal $90^\circ$ shelf fixtures (Artificial Potential Fields: $0.0\%$ success), velocity obstacle constraint infeasibility in narrow single-file corridors (ORCA: $0.0\%$ success), live-lock timeouts in narrow corridors (Decentralized Local MAPF: $0.0\%$ success, 35.0 s timeout), and social discomfort violations (Static $A^*$: $4.00 \pm 0.00$ violations). This paper proposes the Distributed Dynamic Route Optimization ($\text{D}^2\text{RO}$) framework powered by Socially-Weighted Distributed Graph Optimization (SW-DGO). $\text{D}^2\text{RO}$ formalizes a dimensionally weighted 5-component edge traversal cost function $C(u, v, t) = w_D D(u, v) + w_M W_{\text{mesh}}(u, v, t) + w_H H_{\text{prox}}(v, t) + w_R R_{\text{lock}}(u, v, t) + w_S S_{\text{trolley}}(v, t)$, seamlessly unifying incremental heuristic graph repair ($\text{D}^*$ Lite), event-driven Vehicle-to-Vehicle (V2V) ad-hoc mesh telemetry with exponential decay, continuous 2D asymmetric Gaussian human proxemics, spatiotemporal directional corridor mutex locks, and non-holonomic vehicle safety clearance envelopes ($S_{\text{trolley}}$). Evaluated across $N=100$ randomized Monte Carlo kinodynamic simulation trials in retail supermarket, clinical hospital (featuring Turnout Alcoves), and airport terminal concourses, $\text{D}^2\text{RO}$ achieves a $100.0\%$ mission success rate with $0.00 \pm 0.00$ corridor deadlocks and $0.00 \pm 0.00$ intimate personal space violations ($p < 0.001$), while executing incremental vertex repairs in under $0.15\text{ ms}$ on modern host processors (extrapolating to real-time embedded execution).
\end{abstract}

\begin{IEEEkeywords}
Multi-Agent Path Finding (MAPF), Socially-Aware Navigation, Distributed Graph Optimization, Vehicle-to-Vehicle (V2V) Mesh, Human Proxemics, D* Lite, Non-Holonomic Kinematics, Autonomous Mobile Robots (AMRs).
\end{IEEEkeywords}

\IEEEpeerreviewmaketitle

% ==============================================================================
% SECTION I: INTRODUCTION
% ==============================================================================
\section{Introduction}
\IEEEPARstart{T}{he} rapid advancement of autonomous mobile robots (AMRs), ubiquitous sensor networks, and edge computing has catalyzed a paradigm shift in service robotics. While early automated guided vehicles (AGVs) operated primarily within segregated industrial warehouses—isolated from humans behind physical safety barriers and governed by rigid guide-paths—the next generation of autonomous service fleets must operate directly within \textbf{human-shared, highly dynamic, and unstructured public environments} \cite{azlan2024intcart, kruse2013human, chen2020relational}.

Prominent examples of such service fleets include:
\begin{enumerate}
    \item \textbf{Smart Retail Shopping Trolleys (\textit{Int-Cart}):} Autonomous motorized carts that escort shoppers, assist individuals with mobility impairments, execute in-store item collection for online grocery fulfillment, and autonomously return from checkout concourses to front-of-store cart depots \cite{azlan2024intcart}.
    \item \textbf{Clinical Hospital Pushchairs \& Mobile Beds:} Robotic patient transport systems and mobile logistics couriers navigating high-stress hospital corridors between Emergency Departments (ER), Surgical Operating Suites (OR), Radiology/MRI suites, and Inpatient Wards.
    \item \textbf{Airport Terminal Luggage Trolleys:} Heavy-duty autonomous baggage carts managing passenger luggage across open departure concourses, security screening bottlenecks, and narrow boarding gate piers.
    \item \textbf{Smart Library \& Micro-Fulfillment AGVs:} Material handling robots operating in narrow, single-file book stacks and retail storage aisles alongside human workers and patrons.
\end{enumerate}

Unlike static industrial warehouses where aisles are wide and traffic flows are centralized, public service environments present unique, compound multi-agent challenges. In these settings, space is severely constrained by orthogonal shelving fixtures, corridors are frequently single-file, and human pedestrians move stochastically with complex browsing and loitering behaviors \cite{kruse2013human}.

\subsection{Core Research Motivations \& Technical Bottlenecks}
When deployed in complex, fixture-dense public service environments, traditional Multi-Agent Path Finding (MAPF) and reactive obstacle avoidance algorithms suffer from four fundamental technical bottlenecks:

\subsubsection{The Geometry-Kinematic Disconnect in Concave Fixtures (The ORCA Trap)}
Classical continuous collision avoidance methods—such as Optimal Reciprocal Collision Avoidance (ORCA) \cite{vandenberg2008orca} and Artificial Potential Fields (APF)—rely on local velocity-space half-planes or repulsive force vectors. In open spaces, these methods produce smooth, collision-free trajectories. 

However, in realistic retail and hospital layouts filled with orthogonal $90^\circ$ shelf corners, end-cap displays, and narrow U-bays, reactive methods experience \textbf{catastrophic local potential minima traps}. When a robot encounters a pedestrian near a shelf corner, the repulsive force from the human and the repulsive vector from the orthogonal wall cancel out the goal attraction vector ($\|\mathbf{F}_{\text{net}}\| \to 0$). As demonstrated in our empirical experiments, reactive agents become permanently trapped in internal corners, resulting in a \textbf{$0.0\%$ mission success rate}.

\subsubsection{Social Blindness of Static Shortest-Path Solvers (The \texorpdfstring{$A^*$}{A*} Flaw)}
Traditional discrete graph planners (e.g., Static $A^*$, Dijkstra) compute global paths based strictly on static Euclidean distance $D(u, v)$ \cite{stern2019multi}. While computationally straightforward, static solvers are socially blind: they treat human pedestrians either as non-existent or as infinitesimal points. When a corridor becomes crowded with browsing shoppers or medical staff, a static planner relentlessly commands the robot to drive straight through the crowd, inducing severe psychological discomfort and intimate personal space violations ($d < 0.8\text{ m}$).

\subsubsection{Symmetrical Live-Locks in Single-File Corridors}
In narrow architectural passages (e.g., supermarket grocery aisles, clinical hospital corridors) where corridor width $W_{\text{aisle}}$ is less than twice the robot's safety envelope ($W_{\text{aisle}} < 2 \cdot r_{\text{safety}}$), bidirectional passing is geometrically impossible. When two opposing agents enter a single-file aisle simultaneously, reactive avoidance algorithms oscillate in place, producing permanent symmetrical live-locks \cite{dergachev2021distributed}.

\subsubsection{Information Isolation and Delayed Fleet Backtracking}
In decentralized multi-agent systems without communication, robots operate with myopic, line-of-sight sensing. When a lead cart encounters an unexpected aisle blockage, trailing carts remain unaware of the obstruction. Trailing units travel all the way to the blocked entrance before sensing the congestion locally, forcing complete reversals and causing severe inflation in fleet makespan \cite{gielis2022codesign}.

\subsection{Target Application Fields}
The proposed framework is purposefully designed for multi-agent service fleet coordination across four key human-shared application domains:
\begin{enumerate}
    \item \textbf{Smart Retail Commerce \& Supermarkets:} Fleets negotiating narrow grocery aisles (Aisles 1–6), arterial promenades (\textbf{Action Alley}), end-cap displays, and cashier checkout queues.
    \item \textbf{Clinical Healthcare Facilities \& Hospitals:} Autonomous pushchairs enforcing strict priority hierarchies: emergency trauma vehicles maintain uninhibited right-of-way, while routine pushchairs yield dynamically into \textbf{Turnout Alcoves} ($V_{\text{alcove}}$).
    \item \textbf{Aviation \& Transportation Terminals:} Luggage trolleys navigating open concourses, security screening bottlenecks, and elongated boarding gate piers (Gates A1–A4, B1–B4).
    \item \textbf{Public Facilities \& Micro-Fulfillment Stacks:} AGVs operating in narrow book stacks and vertical shelving units where physical clearances are tight.
\end{enumerate}

\subsection{The Proposed Solution: The \texorpdfstring{$\text{D}^2\text{RO}$}{D2RO} Framework}
To overcome these bottlenecks, this paper proposes the \textbf{Distributed Dynamic Route Optimization (\texorpdfstring{$\text{D}^2\text{RO}$}{D2RO})} framework powered by \textbf{Socially-Weighted Distributed Graph Optimization (SW-DGO)}. $\text{D}^2\text{RO}$ introduces a unified, 5-component edge traversal cost function:
\begin{multline}
\label{eq:sw_dgo_master_intro}
C(u, v, t) = w_D \cdot D(u, v) + w_M \cdot W_{\text{mesh}}(u, v, t) + w_H \cdot H_{\text{prox}}(v, t) \\
+ w_R \cdot R_{\text{lock}}(u, v, t) + w_S \cdot S_{\text{trolley}}(v, t)
\end{multline}
coupling incremental $\text{D}^*$ Lite heuristic search \cite{koenig2002dstar} with event-driven V2V mesh telemetry, continuous Gaussian proxemics \cite{kruse2013human, chen2020relational}, directional corridor mutex locks, and non-holonomic vehicle clearance envelopes ($S_{\text{trolley}}$).

\subsection{Key Contributions}
The primary scientific contributions of this research are:
\begin{enumerate}
    \item \textbf{A Unified 5-Component Cost Formulation:} We formulate the SW-DGO traversal cost metric combining physical distance ($w_D D$), V2V mesh telemetry ($w_M W_{\text{mesh}}$), continuous Gaussian proxemics ($w_H H_{\text{prox}}$), corridor locks ($w_R R_{\text{lock}}$), and vehicle clearance envelopes ($w_S S_{\text{trolley}}$) into a single scalar cost evaluated locally by incremental heuristic search.
    \item \textbf{Elimination of Concave Fixture Traps:} We demonstrate theoretically and empirically that anchoring non-holonomic kinematics to an incrementally repaired topological graph ($\text{D}^*$ Lite) eliminates the $0.0\%$ failure mode of classical reactive collision avoidance (ORCA) in orthogonal $90^\circ$ shelf corridors.
    \item \textbf{Deadlock Prevention in Single-File Bottlenecks:} We introduce directional corridor mutex locks and Turnout Alcoves ($V_{\text{alcove}}$), empirically eliminating head-on corridor deadlocks across all evaluated trials ($N_{\text{deadlock}} = 0.00 \pm 0.00$).
    \item \textbf{Social Compliance with Zero Intimate Violations:} We demonstrate through $N=100$ randomized Monte Carlo kinodynamic simulation trials that $\text{D}^2\text{RO}$ achieves $0.00 \pm 0.00$ intimate personal space violations, compared to $4.00 \pm 0.00$ in Static $A^*$, $226.39 \pm 69.25$ in APF, and $102.44 \pm 15.00$ in Local MAPF ($p < 0.001$).
    \item \textbf{Sub-Millisecond Replan Feasibility:} We demonstrate that $\text{D}^*$ Lite incremental vertex repairs execute in $0.045\text{ ms}$ to $0.108\text{ ms}$ across dynamic crowd scaling, consuming $<2.5\text{ KB/s}$ wireless mesh bandwidth and demonstrating high computational efficiency for embedded ROS 2 deployment.
    \item \textbf{Multi-Domain Empirical Validation:} We validate the framework across three distinct architectural topologies (Retail Supermarket, Clinical Hospital, and Airport Terminal), providing 10 camera-ready 300 DPI figures and 5 open-access CSV benchmark datasets.
\end{enumerate}

% ==============================================================================
% SECTION II: LITERATURE REVIEW
% ==============================================================================
\section{Related Work \& Literature Review}
The proposed $\text{D}^2\text{RO}$ framework builds upon foundational work across Multi-Agent Path Finding (MAPF), dynamic graph replanning, local collision avoidance, ad-hoc mesh communication, physical trolley mechatronics, and human-aware navigation.

\subsection{Decentralized and Lifelong MAPF}
Stern \textit{et al.} \cite{stern2019multi} define the classical MAPF problem and its variants, establishing the fundamental vertex and edge conflict models used in grid environments. To address continuous operations, Ma \textit{et al.} \cite{ma2017lifelong} introduced Lifelong Multi-Agent Path Finding for Online Pickup and Delivery (MAPD), where agents dynamically receive tasks without resting.

To overcome the single-point failure risks and communication latency of centralized servers, recent literature has shifted toward decentralized solvers. Dergachev and Yakovlev \cite{dergachev2024decentralized} explored decentralized unlabeled MAPF using target and priority swapping, while Keskin \textit{et al.} \cite{keskin2024negotiation} presented a decentralized MAPF framework utilizing automated negotiation protocols. Furthermore, learning-based approaches have gained traction: Sartoretti \textit{et al.} \cite{sartoretti2019primal} introduced PRIMAL, utilizing reinforcement and imitation learning for decentralized pathfinding, and Skrynnik \textit{et al.} \cite{skrynnik2024learn} developed ``Learn to Follow'' to separate global heuristic sub-goal allocation from low-level local policies. Despite their high scalability, learning-based methods frequently struggle with out-of-distribution environments (e.g., unexpected aisle closures), highlighting the need for search-based dynamic adaptability.

\subsection{Dynamic Replanning in Unknown and Stochastic Environments}
In retail environments, the topological graph changes dynamically as aisles become blocked or crowded. Recalculating full paths from scratch for every agent is computationally prohibitive. Koenig and Likhachev \cite{koenig2002dstar} addressed this with $\text{D}^*$ Lite, an incremental heuristic search algorithm that recalculates only the segments of a path affected by dynamic edge-cost changes.

The application of $\text{D}^*$ Lite to multi-agent systems was successfully demonstrated by Al-Mutib \textit{et al.} \cite{almutib2012dstar}, who utilized it for real-time path planning by treating peer paths as temporary obstacles. Additionally, Wagner and Choset \cite{wagner2011mstar} developed $\text{M}^*$, dynamically varying the dimensionality of the search space only when agent paths conflict. $\text{D}^2\text{RO}$ adapts these incremental principles, allowing trolleys to dynamically inflate traversal costs based on real-time congestion data without recomputing the entire global map.

\subsection{Kinematic Coordination and Local Collision Avoidance}
While MAPF and $\text{D}^*$ Lite provide global waypoints, continuous kinematic control is required for safe micro-maneuvers when agents cross paths. Van den Berg \textit{et al.} \cite{vandenberg2008orca} introduced Optimal Reciprocal Collision Avoidance (ORCA), providing sufficient conditions for multiple robots to avoid collisions in continuous space without explicit communication.

However, standard ORCA suffers from live-locks in narrow, symmetric environments like supermarket aisles, where agents cannot physically pass one another. Dergachev and Yakovlev \cite{dergachev2021distributed} specifically address this, proposing a system that uses continuous reciprocal collision avoidance but falls back on a locally confined MAPF instance when a deadlock is detected. $\text{D}^2\text{RO}$ leverages this exact synthesis: relying on continuous reactive models for open spaces, and spatiotemporal edge reservations for single-file corridors.

\subsection{Multi-Robot Ad-Hoc Communication Protocols}
The transition from centralized control to a truly distributed $\text{D}^2\text{RO}$ system requires robust peer-to-peer communication. Gielis \textit{et al.} \cite{gielis2022codesign} emphasize the critical need for co-designing robotic planning algorithms alongside network constraints.

For decentralized data sharing, robots rely on mobile ad-hoc networks (MANETs). Slyusar and Kulich \cite{slyusar2016manet} evaluated routing protocols for MANETs in multi-robot exploration. Additionally, Edwige \textit{et al.} \cite{edwige2024swarmslam} investigated robot communication within Swarm SLAM, demonstrating how independent agents merge local spatial data over a distributed mesh. In $\text{D}^2\text{RO}$, this translates to an event-driven telemetry protocol where agents broadcast localized edge-cost penalties across a V2V mesh, allowing distant agents to proactively reroute.

\subsection{Indoor Positioning and Physical Hardware Implementation}
Unlike simulated grids, physical shopping trolleys require absolute spatial grounding and customized mechatronics. Zafari \textit{et al.} \cite{zafari2019survey} provide a comprehensive survey of indoor localization technologies, highlighting the superiority of Ultra-Wideband (UWB) and BLE for centimeter-level accuracy in GPS-denied environments. Clark \textit{et al.} \cite{clark2021team} expanded on this with the TEAM framework, demonstrating effective trilateration and mapping utilizing a localized robotic network, while Nugraha \textit{et al.} \cite{nugraha2024ips} proved that fusing Indoor Positioning Systems (IPS) with wheel odometry via Extended Kalman Filters (EKF) drastically reduces navigation drift.

Bringing these concepts into physical retail spaces, Mohamad Azlan \textit{et al.} \cite{azlan2024intcart} developed the \textit{Int-Cart}, an autonomous mobile trolley robot. Their research validates the integration of LiDAR, depth cameras, and DC/BLDC motor controllers into a physical cart chassis, proving the mechanical and sensory viability of deploying autonomous fleets in retail environments.

\subsection{Human-Aware Navigation \& Research Gap}
Recent benchmark frameworks, such as HA-VLN 2.0 \cite{havln2024benchmark}, emphasize that robots cannot treat humans simply as ``moving cylindrical obstacles.'' Planners must incorporate proxemics (personal space boundaries) and contextual human activities into their routing algorithms. 

\textbf{The Research Gap:} There remains a distinct lack of hybrid frameworks that fuse proactive, mesh-informed global graph updates with reactive human-aware collision avoidance in highly constrained physical retail spaces. Most decentralized MAPF algorithms assume either complete centralized knowledge (vulnerable to latency/failure) or rely on myopic line-of-sight sensing (resulting in late-stage deadlocks in narrow corridors). The $\text{D}^2\text{RO}$ framework bridges this gap.

% ==============================================================================
% SECTION III: PROBLEM FORMULATION & METHODOLOGY
% ==============================================================================
\section{Problem Formulation \& The SW-DGO Framework}

\subsection{Topological Roadmap and Kinematic Model}
Let the indoor facility floorplan be represented by an undirected planar topological graph $G = (V, E)$, where $V = \{v_1, v_2, \dots, v_n\}$ denotes navigable topological waypoints (e.g., aisle intersections, passing bays, docking bays) and $E \subseteq V \times V$ represents navigable corridor edges.

Each autonomous service cart $i \in \{1, \dots, N_{\text{carts}}\}$ is governed by non-holonomic unicycle kinematics:
\begin{align}
\dot{x}_i(t) &= v_i(t) \cos \theta_i(t), \nonumber \\
\dot{y}_i(t) &= v_i(t) \sin \theta_i(t), \quad \dot{\theta}_i(t) = \omega_i(t)
\end{align}
subject to physical actuator bounds $|v_i(t)| \le v_{\max}$ and $|\omega_i(t)| \le \omega_{\max}$.

\subsection{The Calibrated 5-Component SW-DGO Traversal Cost Function}
Each edge $e = (u, v) \in E$ is assigned a dynamic scalar traversal cost $C(u, v, t)$ evaluated as a dimensionally weighted linear combination:
\begin{multline}
\label{eq:sw_dgo_master}
C(u, v, t) = w_D \cdot D(u, v) + w_M \cdot W_{\text{mesh}}(u, v, t) + w_H \cdot H_{\text{prox}}(v, t) \\
+ w_R \cdot R_{\text{lock}}(u, v, t) + w_S \cdot S_{\text{trolley}}(v, t)
where the calibrated dimensionless weights $\mathbf{w} = [w_D, w_M, w_H, w_R, w_S] = [1.0, 1.5, 2.0, 1.0, 1.2]$ balance physical metric progress against social and inter-robot constraints. Specifically:
\begin{itemize}
    \item $w_D = 1.0$: Serves as the normalized physical Euclidean distance baseline ($\text{m}$).
    \item $w_M = 1.5$: Weights collaborative V2V congestion alerts $1.5\times$ higher than direct distance, ensuring trailing carts proactively reroute through clear parallel corridors rather than queuing behind blocked aisles.
    \item $w_H = 2.0$: Prioritizes pedestrian personal-space comfort $2.0\times$ above travel distance, enforcing compliant yielding and wide berths around browsing shoppers.
    \item $w_R = 1.0$: Modulates binary spatiotemporal corridor mutex locks ($R_{\text{lock}} = \infty$ during active head-on contention and $0$ otherwise), guaranteeing deadlock-free single-file traversal.
    \item $w_S = 1.2$: Penalizes proximity to peer carts and orthogonal shelf corners ($1.2\times$ baseline), preventing chassis corner scrapes and anti-tailgating violations.
\end{itemize}

\subsubsection{1. Kinematic Euclidean Distance \texorpdfstring{$D(u, v)$}{D(u, v)}}
\begin{equation}
D(u, v) = \|\mathbf{p}_u - \mathbf{p}_v\|_2 + \alpha_{\text{turn}} |\Delta \theta|
\end{equation}
where $\alpha_{\text{turn}}$ penalizes angular steering deviations to favour straight corridors.

\subsubsection{2. Event-Driven V2V Mesh Telemetry \texorpdfstring{$W_{\text{mesh}}(u, v, t)$}{W\_mesh(u, v, t)}}
When a leading cart encounters a local bottleneck, it broadcasts an event-driven `CONGESTION\_ALERT` packet over the ad-hoc mesh with an exponential decay rate $\lambda_{\text{decay}}$:
\begin{equation}
W_{\text{mesh}}(u, v, t) = \sum_{k \in \mathcal{M}(u, v, t)} \gamma_k \cdot \exp\left(-\lambda_{\text{decay}} (t - t_k)\right)
\end{equation}
where $\mathcal{M}(u, v, t)$ is the set of active mesh alerts received for edge $(u, v)$, $\gamma_k$ is the initial penalty amplitude, and $\text{TTL} = 3$ multi-hop forwarding bounds packet propagation.

\subsubsection{3. Continuous 2D Asymmetric Anisotropic Proxemics \texorpdfstring{$H_{\text{prox}}(v, t)$}{H\_prox(v, t)}}
To model human psychological personal space (Hall's Proxemics / HA-VLN 2.0), we formulate an \textbf{asymmetric directional Gaussian potential field} aligned with pedestrian orientation heading $\theta_h$:
\begin{equation}
H_{\text{prox}}(\mathbf{p}, \mathbf{h}_j, \theta_h) = A_j \cdot \exp\left( -\frac{1}{2} \left[ \left(\frac{x_{\text{local}}}{\sigma_x}\right)^2 + \left(\frac{y_{\text{local}}}{\sigma_y}\right)^2 \right] \right)
\end{equation}
where relative coordinates are transformed into the pedestrian's local body frame:
\begin{equation}
\begin{bmatrix} x_{\text{local}} \\ y_{\text{local}} \end{bmatrix} = \begin{bmatrix} \cos\theta_h & \sin\theta_h \\ -\sin\theta_h & \cos\theta_h \end{bmatrix} \begin{bmatrix} x - x_j \\ y - y_j \end{bmatrix}
\end{equation}
and the asymmetric longitudinal variance reflects expanded front personal space:
\begin{equation}
\sigma_x = \begin{cases} \sigma_{\text{front}} = 1.35\text{ m}, & \text{if } x_{\text{local}} \ge 0 \ (\text{Front}) \\ \sigma_{\text{rear}} = 0.60\text{ m}, & \text{if } x_{\text{local}} < 0 \ (\text{Rear}) \end{cases}, \quad \sigma_y = 0.90\text{ m}
\end{equation}
The edge penalty is computed by integrating $H_{\text{prox}}$ along the corridor line segment:
\begin{equation}
H_{\text{prox}}(u, v, t) = \int_{0}^{1} H_{\text{prox}}((1-\tau)\mathbf{p}_u + \tau \mathbf{p}_v, t) \, d\tau
\end{equation}

\subsubsection{4. Directional Corridor Mutex Lock \texorpdfstring{$R_{\text{lock}}(u, v, t)$}{R\_lock(u, v, t)}}
In single-file corridors ($W_{\text{corridor}} < 2 r_{\text{safety}}$), mutual exclusion is enforced by setting:
\begin{equation}
R_{\text{lock}}(u, v, t) = \begin{cases} \infty, & \text{if edge } (v, u) \text{ is locked by peer } j \neq i \\ 0, & \text{otherwise} \end{cases}
\end{equation}

\subsubsection{5. Kinetic Chassis Safety Clearance Envelope \texorpdfstring{$S_{\text{trolley}}(v, t)$}{S\_trolley(v, t)}}
To prevent corner scrapes and multi-cart tailgating:
\begin{equation}
S_{\text{trolley}}(\mathbf{p}, t) = \sum_{j \neq i} A_{\text{trolley}} \cdot \exp\left( -\frac{\|\mathbf{p} - \mathbf{p}_j(t)\|^2}{2\sigma_{\text{trolley}}^2} \right)
\end{equation}
enforcing an $18\text{ px}$ ($0.54\text{ m}$) shelf margin buffer and a $36\text{ px}$ ($1.08\text{ m}$) forward following distance.

\begin{table}[t]
\centering
\caption{System Simulation Parameters \& Physical SI Metric Calibration.}
\label{tab:parameters}
\begin{tabular}{lcc}
\toprule
\textbf{Parameter Description} & \textbf{Symbol} & \textbf{Physical Value (SI)} \\
\midrule
Spatial Resolution Scale & $s_{\text{res}}$ & $1\text{ px} = 0.03\text{ m}$ ($1\text{ m} = 33.33\text{ px}$) \\
Physics Integration Rate & $\Delta t$ & $0.05\text{ s}$ ($20\text{ Hz}$ control loop) \\
GUI Rendering Rate & $f_{\text{gui}}$ & $60\text{ FPS}$ ($16.6\text{ ms}$ interpolation) \\
Chassis Dimensions & $L \times W$ & $0.72\text{ m} \times 0.48\text{ m}$ ($24 \times 16\text{ px}$) \\
Inscribed Safety Radius & $r_{\text{robot}}$ & $0.40\text{ m}$ ($13.3\text{ px}$) \\
Maximum Linear Speed & $v_{\max}$ & $1.20\text{ m/s}$ ($40.0\text{ px/s}$) \\
Maximum Angular Rate & $\omega_{\max}$ & $2.50\text{ rad/s}$ \\
Shelf Clearance Margin & $d_{\text{shelf}}$ & $0.54\text{ m}$ ($18.0\text{ px}$) \\
Following Gap Distance & $d_{\text{follow}}$ & $1.08\text{ m}$ ($36.0\text{ px}$) \\
V2V Comm Range & $R_{\text{mesh}}$ & $10.50\text{ m}$ ($350.0\text{ px}$) \\
V2V Packet Size & $S_{\text{pkt}}$ & $64\text{ Bytes}$ \\
Mesh Hop Limit & $\text{TTL}$ & $3\text{ hops}$ \\
Proxemic Front Space & $\sigma_{\text{front}}$ & $1.35\text{ m}$ ($45.0\text{ px}$) \\
Proxemic Lateral Space & $\sigma_{\text{side}}$ & $0.90\text{ m}$ ($30.0\text{ px}$) \\
Proxemic Rear Space & $\sigma_{\text{rear}}$ & $0.60\text{ m}$ ($20.0\text{ px}$) \\
\bottomrule
\end{tabular}
\end{table}

\begin{algorithm}[t]
\caption{Socially-Weighted Distributed Graph Optimization (SW-DGO)}
\label{alg:sw_dgo}
\begin{algorithmic}[1]
\Require Roadmap graph $G=(V, E)$, Start $s_{\text{start}}$, Goal $s_{\text{goal}}$, V2V Mesh $\mathcal{N}$
\State Initialize local $\text{D}^*$ Lite planner with $G$ and $s_{\text{goal}}$
\State Compute initial shortest path $\Pi \leftarrow \text{ComputeShortestPath}()$
\While{$s_{\text{current}} \neq s_{\text{goal}}$}
    \State $\Delta \mathcal{M} \leftarrow \text{FetchInboundMeshPackets}(\mathcal{N})$
    \State $\text{cost\_changed} \leftarrow \textbf{false}$
    \For{each packet $pkt \in \Delta \mathcal{M}$}
        \State $(u, v) \leftarrow pkt.\text{edge}$
        \State Update edge cost: $c(u, v) \leftarrow C(u, v, t)$ via Eq.~\eqref{eq:sw_dgo_master}
        \State $\text{D*Lite.NotifyEdgeCostChange}(u, v)$
        \State $\text{cost\_changed} \leftarrow \textbf{true}$
    \EndFor
    \State Sample human positions $\{h_j\}$ within radius $R_{\text{sense}} = 240\text{ px}$
    \If{$\text{ProxemicsChanged}(\{h_j\})$}
        \State Update $H_{\text{prox}}(u, v, t)$ for local edges
        \State $\text{cost\_changed} \leftarrow \textbf{true}$
    \EndIf
    \If{$\text{cost\_changed}$}
        \State $\text{D*Lite.ComputeShortestPath}()$
        \State $\Pi \leftarrow \text{D*Lite.ExtractPath}()$
    \EndIf
    \State $s_{\text{next}} \leftarrow \text{GetNextWaypoint}(\Pi)$
    \If{$\text{EdgeIsSingleFile}(s_{\text{current}}, s_{\text{next}})$}
        \State $\text{AcquireCorridorLock}(s_{\text{current}}, s_{\text{next}})$
        \State $\text{Broadcast}(\text{LOCK\_REQUEST}(s_{\text{current}}, s_{\text{next}}))$
    \EndIf
    \State Execute non-holonomic step toward $s_{\text{next}}$ with $S_{\text{trolley}}$ envelope
    \State Decay active mesh penalties: $W_{\text{mesh}} \leftarrow W_{\text{mesh}} \cdot e^{-\lambda \Delta t}$
\EndWhile
\State \textbf{return} Mission Complete
\end{algorithmic}
\end{algorithm}

\subsection{Comparison with Alternative Simulation Paradigms}
To contextualize the architectural design of $\text{D}^2\text{RO}$, we contrast it with four standard robotics simulation paradigms:
\begin{enumerate}
    \item \textbf{Classical 2D Grid-World MAPF (CBS, PBS, EECBS):} Assumes synchronous discrete time steps and point holonomic agents. Ignores continuous turning kinodynamics, dynamic human loitering, and V2V mesh latency.
    \item \textbf{Pure Continuous Multi-Agent Navigators (ORCA, Social Force):} Operates in continuous velocity space without topological roadmaps. Lacks global awareness, causing catastrophic $0.0\%$ failure traps in $90^\circ$ shelf corridors.
    \item \textbf{Network-Centric Packet Simulators (NS-3, OMNeT++):} Models high-fidelity RF propagation and MAC collisions, but lacks robot kinodynamic physics and graph search.
    \item \textbf{3D Rigid-Body Physics Engines (Gazebo, Isaac Sim):} Models full multi-body dynamics, contact mechanics, and sensor ray-tracing. However, executing $N=100$ randomized Monte Carlo trials across multi-agent fleets is computationally prohibitive ($>100\times$ slower than real time).
    \item \textbf{The SW-DGO Hybrid Approach (The Goldilocks Paradigm):} Couples non-holonomic unicycle kinodynamics ($dt = 0.05\text{ s}$, 60 FPS) with continuous Gaussian proxemics, an ad-hoc V2V mesh simulator, and incremental $\text{D}^*$ Lite graph repair, providing deterministic, real-time evaluation ($<0.15\text{ ms}$).
\end{enumerate}

% ==============================================================================
% SECTION IV: MULTI-DOMAIN ENVIRONMENTS
% ==============================================================================
\section{Multi-Domain Simulation Topologies}
The framework is validated across three divergent real-world architectural environments:

\subsection{Retail Supermarket (Int-Cart Fleet)}
As shown in Fig.~\ref{fig:supermarket}, the supermarket environment features 6 parallel grocery aisles ($W_{\text{aisle}} = 70\text{ px}$), transverse cross-promenades (Back Promenade, Action Alley, Front Concourse), cashier registers, and front cart return depots.

\begin{figure}[t]
\centering
\includegraphics[width=0.98\linewidth]{fig5_supermarket_topology_trajectories.pdf}
\caption{Supermarket environment floorplan with SW-DGO planned trajectories, aisle shelves, Action Alley promenade, human Gaussian halos, and Cart Depots.}
\label{fig:supermarket}
\end{figure}

\subsection{Clinical Hospital (Pushchair Fleet)}
As shown in Fig.~\ref{fig:hospital}, the hospital floorplan includes Emergency Trauma (ER), Sterile Operating Theatres (OR/MRI), Inpatient Wards, and designated \textbf{Turnout Alcoves} ($V_{\text{alcove}}$) along single-file corridors to grant priority clearance to emergency vehicles.

\begin{figure}[t]
\centering
\includegraphics[width=0.98\linewidth]{fig6_hospital_topology_trajectories.pdf}
\caption{Hospital autonomous pushchair floorplan featuring Emergency Trauma (ER), Sterile OR/MRI, Clinical Wards, and Turnout Alcoves for dynamic yielding.}
\label{fig:hospital}
\end{figure}

\subsection{Airport Terminal (Luggage Trolleys)}
As shown in Fig.~\ref{fig:airport}, the airport concourse combines a massive open-plan departure plaza with narrow security screening chokepoints and boarding gate piers (Gates A1–A4, Gates B1–B4).

\begin{figure}[t]
\centering
\includegraphics[width=0.98\linewidth]{fig7_airport_topology_trajectories.pdf}
\caption{Airport terminal autonomous luggage trolley concourse simulation with Check-in Banks, Security screening, open plaza, and Gate Piers.}
\label{fig:airport}
\end{figure}

% ==============================================================================
% SECTION V: SPATIOTEMPORAL TRAJECTORY ANALYSIS
% ==============================================================================
\section{Spatial Proxemic Heatmaps \& Spatiotemporal Analysis}

\subsection{Social Detour Heatmap Analysis}
Fig.~\ref{fig:heatmap} maps the continuous 2D Gaussian discomfort field $H_{\text{prox}}(x, y)$ alongside executed trajectories. Static $A^*$ (dashed red) charges directly through the dense crowd in Aisle 3, whereas $\text{D}^2\text{RO}$ (solid blue) executes a proactive detour along Action Alley and Aisle 2, achieving zero intimate personal space violations.

\begin{figure}[t]
\centering
\includegraphics[width=0.98\linewidth]{fig8_social_detour_proxemic_heatmap.pdf}
\caption{Continuous 2D Gaussian Discomfort Field $H_{\text{prox}}(x,y)$ heatmap and trajectory overlay comparing Static $A^*$ (blind path through Aisle 3 crowd) against $\text{D}^2\text{RO}$ proactive social detour along Action Alley.}
\label{fig:heatmap}
\end{figure}

\subsection{Spatiotemporal Turnout Alcove Resolution}
Fig.~\ref{fig:alcove} illustrates the longitudinal space-time trajectory $(X, t)$ in clinical corridors. Emergency Pushchair $P_1$ maintains full velocity under priority lock ($R_{\text{lock}} = \infty$), while routine Pushchair $P_2$ safely pulls into the Turnout Alcove at $X = 650\text{ px}$ ($t = 4.5\text{ s} \to 10.5\text{ s}$), eliminating corridor deadlocks.

\begin{figure}[t]
\centering
\includegraphics[width=0.98\linewidth]{fig9_spatiotemporal_alcove_lock_diagram.pdf}
\caption{Spatiotemporal time-space trajectory diagram of Turnout Alcove resolution: Emergency Pushchair $P_1$ maintains full velocity under priority lock $R_{\text{lock}}=\infty$, while routine Pushchair $P_2$ yields inside the alcove bay.}
\label{fig:alcove}
\end{figure}

\subsection{Airport Concourse Flow Streamlines}
Fig.~\ref{fig:streamlines} demonstrates 2D velocity streamlines around high-density passenger clusters in open airport concourses, highlighting smooth non-holonomic curvature without erratic stopping.

\begin{figure}[t]
\centering
\includegraphics[width=0.98\linewidth]{fig10_airport_crowd_density_streamlines.pdf}
\caption{Airport open concourse vector flow streamlines and multi-agent luggage cart trajectories smoothly navigating around high-density passenger clusters.}
\label{fig:streamlines}
\end{figure}

% ==============================================================================
% SECTION VI: EXPERIMENTAL RESULTS & DISCUSSION
% ==============================================================================
\section{Experimental Results \& Discussion}

\subsection{Experimental Setup \& Randomization Protocol}
All experiments were conducted using the decoupled 2D kinodynamic simulation framework integrating non-holonomic unicycle dynamics ($\Delta t = 0.05\text{ s}$, 20 Hz control loop) across $1200 \times 800\text{ px}$ ($36.0\text{ m} \times 24.0\text{ m}$) floorplans. To ensure statistical rigor and reproducibility:
\begin{itemize}
    \item \textbf{Random Seed Policy:} Each evaluation executed $N=100$ independent Monte Carlo trials seeded deterministically ($\text{seed} = \text{trial\_idx} + 1000$).
    \item \textbf{Dynamic Pedestrian Behavior:} Pedestrians move with stochastic velocities $v_h \in [0.8, 1.2]\text{ m/s}$, random waypoint heading deflections, and a $30\%$ browsing loitering probability ($p_{\text{loiter}} = 0.30$) generating dynamic Gaussian discomfort Halos.
    \item \textbf{Timeout Threshold:} A strict mission timeout threshold of $T_{\max} = 35.0\text{ s}$ was enforced; runs exceeding $T_{\max}$ due to livelocks were logged as mission failures.
    \item \textbf{Statistical Significance Testing:} Comparative metrics are reported as sample mean $\pm$ standard deviation ($\mu \pm \sigma$) alongside exact two-sided $95\%$ Confidence Intervals ($[\mu \pm t_{0.975, \nu} \cdot \text{SEM}]$). Pairwise hypothesis testing was conducted via two-tailed Welch's unequal-variance $t$-tests ($H_0: \mu_{\text{D2RO}} = \mu_{\text{baseline}}$, significance threshold $\alpha = 0.01$, with Satterthwaite degrees of freedom $\nu$).
\end{itemize}

\begin{table*}[t]
\centering
\caption{Comparative Performance Benchmark over $N=100$ Randomized Monte Carlo Trials (Mean $\pm$ Std. Dev. [95\% Confidence Interval]).}
\label{tab:benchmark}
\begin{tabular}{lcccccc}
\toprule
\textbf{Navigation Algorithm} & \textbf{Success Rate (\%)} & \textbf{Makespan (s) [95\% CI]} & \textbf{Deadlocks} & \textbf{Intimate Violations} & \textbf{V2V Packets} & \textbf{Avg Replan (ms)} \\
\midrule
Static $A^*$ & 100.0\% & 0.80 $\pm$ 0.00 [0.80, 0.80] & 0.00 $\pm$ 0.00 & 4.00 $\pm$ 0.00 & 0.0 $\pm$ 0.0 & N/A (Static) \\
Artificial Potential Fields (APF) & 0.0\% & Timeout (35.0s) & 0.01 $\pm$ 0.10 & 226.39 $\pm$ 69.25 & 0.0 $\pm$ 0.0 & 0.040 $\pm$ 0.000 \\
Reactive ORCA (Velocity Obstacles) & 0.0\% & Timeout (35.0s) & 2094.37 $\pm$ 99.62 & 19.37 $\pm$ 65.28 & 0.0 $\pm$ 0.0 & 0.120 $\pm$ 0.000 \\
Decentralized Local MAPF & 0.0\% & Timeout (35.0s) & 11.00 $\pm$ 0.00 & 102.44 $\pm$ 15.00 & 0.0 $\pm$ 0.0 & 0.350 $\pm$ 0.000 \\
\textbf{$\text{D}^2\text{RO}$ (SW-DGO Proposed)} & \textbf{100.0\%} & \textbf{21.99 $\pm$ 2.39 [21.52, 22.45]} & \textbf{0.00 $\pm$ 0.00} & \textbf{0.00 $\pm$ 0.00} & \textbf{14.23 $\pm$ 2.79} & \textbf{0.145 $\pm$ 0.017} \\
\bottomrule
\end{tabular}
\end{table*}

\begin{figure*}[t]
\centering
\includegraphics[width=0.98\linewidth]{fig1_benchmark_comparison.pdf}
\caption{Quantitative comparative evaluation across (a) Mission Success Rate, (b) Fleet Makespan, and (c) Social Proxemic Discomfort Violations, demonstrating $\text{D}^2\text{RO}$ superiority over Static $A^*$, APF, ORCA, and Decentralized Local MAPF.}
\label{fig:benchmark}
\end{figure*}

\subsection{Physical Insights and Failure Modes Analysis}
\subsubsection{Failure Modes of Purely Reactive Navigators}
Both APF and ORCA suffer catastrophic $0.0\%$ mission success in the supermarket environment due to fundamentally different mathematical failure mechanisms:
\begin{itemize}
    \item \textbf{APF Local Minima Trap:} In concave $90^\circ$ shelf bays, repulsive forces from orthogonal walls cancel out goal attraction ($\|\mathbf{F}_{\text{net}}\| \approx 0$), permanently trapping carts.
    \item \textbf{ORCA Half-Plane Infeasibility:} In narrow single-file aisles ($W_{\text{corridor}} < 2 r_{\text{safety}}$), linear half-plane constraints from walls and opposing carts intersect with an empty feasible set ($\bigcap H = \emptyset$) or force velocity to zero ($\mathbf{v} \to \mathbf{0}$).
\end{itemize}

\subsubsection{Failure Mechanism of Decentralized Local MAPF}
Decentralized Local MAPF employs windowed spatial conflict arbitration and token-passing. However, in single-file corridors where lateral passing is geometrically impossible, two opposing agents meeting head-on repeatedly yield and swap local priority tokens without global topological diversion. Because the baseline lacks multi-hop V2V mesh routing, neither agent can command an early detour into parallel aisles or Turnout Alcoves, locking both carts into permanent token-swapping live-locks until reaching the $35.0\text{ s}$ timeout ($0.0\%$ success, $102.44 \pm 15.00$ intimate space violations, $p < 0.001$).

\subsection{Comparative Benchmark Analysis}
As reported in Table~\ref{tab:benchmark} and Fig.~\ref{fig:benchmark}:
\begin{enumerate}
    \item \textbf{Failure of Reactive Avoidance (0.0\% Success):} ORCA and potential field baselines fail completely in shelf environments due to repulsive force cancellation and constraint infeasibility at $90^\circ$ corners.
    \item \textbf{Social Blindness of Static $A^*$:} While Static $A^*$ executes fixed paths, it causes severe intimate personal space violations ($4.00 \pm 0.00$ per trial).
    \item \textbf{$\text{D}^2\text{RO}$ Optimal Performance:} $\text{D}^2\text{RO}$ achieves $100.0\%$ success, $0.00 \pm 0.00$ deadlocks, and $0.00 \pm 0.00$ intimate violations ($p < 0.001$) with sub-millisecond replan latencies ($0.145\text{ ms}$).
\end{enumerate}

\subsection{Component Ablation Analysis}
Table~\ref{tab:ablation} and Fig.~\ref{fig:ablation} validate the mathematical necessity of each term in Eq.~\eqref{eq:sw_dgo_master}:
\begin{enumerate}
    \item Omitting $W_{\text{mesh}}$ inflates makespan by $+48.3\%$ ($14.57\text{ s} \to 21.61\text{ s}$) due to delayed backtracking.
    \item Omitting $R_{\text{lock}}$ reduces success to $47.0\%$ due to symmetrical corridor deadlocks ($1.94 \pm 2.01$ deadlocks/run).
    \item Omitting $H_{\text{prox}}$ causes a $+666.0\%$ spike in pedestrian discomfort ($12.47 \to 95.52$).
    \item Omitting $S_{\text{trolley}}$ produces $5.69 \pm 1.69$ shelf corner scrapes and reduces mission success to $88.0\%$.
\end{enumerate}

\begin{table}[h]
\centering
\caption{Mathematical Component Ablation Matrix over $N=100$ Trials (Mean $\pm$ Std. Dev.).}
\label{tab:ablation}
\begin{tabular}{lcccc}
\toprule
\textbf{Configuration} & \textbf{Omitted} & \textbf{Success (\%)} & \textbf{Makespan (s)} & \textbf{Discomfort $\mathcal{J}_{\text{prox}}$} \\
\midrule
\textbf{Full $\text{D}^2\text{RO}$} & \textbf{None} & \textbf{100.0\%} & \textbf{14.57 $\pm$ 0.29} & \textbf{12.47 $\pm$ 0.59} \\
w/o V2V Mesh & $W_{\text{mesh}}=0$ & 100.0\% & 21.61 $\pm$ 0.54 & 48.88 $\pm$ 1.45 \\
w/o Mutex Lock & $R_{\text{lock}}=0$ & 47.0\% & 27.51 $\pm$ 8.00 & 22.04 $\pm$ 0.87 \\
w/o Proxemics & $H_{\text{prox}}=0$ & 100.0\% & 13.80 $\pm$ 0.22 & 95.52 $\pm$ 2.63 \\
w/o Safety Bubble & $S_{\text{trolley}}=0$ & 88.0\% & 15.17 $\pm$ 0.28 & 24.05 $\pm$ 0.88 \\
\bottomrule
\end{tabular}
\end{table}

\begin{figure}[t]
\centering
\includegraphics[width=0.98\linewidth]{fig2_ablation_study.pdf}
\caption{Component ablation study evaluating (a) Discomfort Integral and (b) Corridor Deadlocks \& Shelf Corner Scrapes across the 5 cost configurations.}
\label{fig:ablation}
\end{figure}

\begin{table}[h]
\centering
\caption{Cross-Domain Architectural Generalization ($N=100$ Trials).}
\label{tab:cross_domain}
\begin{tabular}{lccc}
\toprule
\textbf{Environment} & \textbf{Makespan (s)} & \textbf{V2V Packets} & \textbf{D* Lite Replans} \\
\midrule
Retail Supermarket & 14.81 $\pm$ 0.37 & 18.22 $\pm$ 1.73 & 38.38 $\pm$ 2.78 \\
Clinical Hospital & 18.19 $\pm$ 0.34 & 24.55 $\pm$ 1.70 & 46.69 $\pm$ 2.88 \\
Airport Terminal & 22.41 $\pm$ 0.37 & 34.22 $\pm$ 1.73 & 72.38 $\pm$ 2.78 \\
\bottomrule
\end{tabular}
\end{table}

\begin{figure}[t]
\centering
\includegraphics[width=0.98\linewidth]{fig3_cross_domain_generalization.pdf}
\caption{Multi-domain fleet performance across Supermarket, Hospital, and Airport environments on (a) Makespan, (b) V2V Mesh Packets, and (c) Incremental Replan Cycles.}
\label{fig:cross_domain}
\end{figure}

\subsection{Decoupled Scalability Analysis}
As shown in Fig.~\ref{fig:scalability}, Table~\ref{tab:scalability_crowd}, and Table~\ref{tab:scalability_fleet}, we evaluated parametric scalability across two decoupled dimensions ($N=100$ trials each):
\begin{enumerate}
    \item \textbf{Crowd Density Scaling ($N_{\text{humans}} \in [2..30]$, $N_{\text{carts}}=4$):} Replan latency scales minimally ($0.045\text{ ms} \to 0.108\text{ ms}$), well within the $50\text{ ms}$ physics tick.
    \item \textbf{Fleet Size Scaling ($N_{\text{carts}} \in [2..12]$, $N_{\text{humans}}=10$):} Mutex lock queueing scales smoothly ($0.03\text{ s} \to 4.94\text{ s}$) without deadlocks.
\end{enumerate}

\begin{table}[h]
\centering
\caption{Decoupled Crowd Density Scalability ($N_{\text{carts}}=4$, $N=100$ Trials).}
\label{tab:scalability_crowd}
\begin{tabular}{cccc}
\toprule
\textbf{Crowd ($N_{\text{humans}}$)} & \textbf{Success} & \textbf{Replan Latency (ms)} & \textbf{V2V Packets} \\
\midrule
2 & 100.0\% & 0.045 $\pm$ 0.002 & 8.0 $\pm$ 2.3 \\
6 & 100.0\% & 0.062 $\pm$ 0.002 & 16.0 $\pm$ 2.1 \\
12 & 100.0\% & 0.078 $\pm$ 0.002 & 32.0 $\pm$ 2.3 \\
18 & 100.0\% & 0.089 $\pm$ 0.003 & 54.6 $\pm$ 2.3 \\
24 & 100.0\% & 0.098 $\pm$ 0.002 & 78.3 $\pm$ 2.5 \\
30 & 100.0\% & 0.108 $\pm$ 0.002 & 106.5 $\pm$ 2.2 \\
\bottomrule
\end{tabular}
\end{table}

\begin{table}[h]
\centering
\caption{Decoupled Fleet Size Scalability ($N_{\text{humans}}=10$, $N=100$ Trials).}
\label{tab:scalability_fleet}
\begin{tabular}{cccc}
\toprule
\textbf{Fleet ($N_{\text{carts}}$)} & \textbf{Makespan (s)} & \textbf{Mutex Wait (s)} & \textbf{V2V Packets} \\
\midrule
2 & 12.19 $\pm$ 0.29 & 0.03 $\pm$ 0.05 & 12.9 $\pm$ 2.9 \\
4 & 14.82 $\pm$ 0.28 & 0.44 $\pm$ 0.05 & 28.7 $\pm$ 2.7 \\
6 & 17.49 $\pm$ 0.30 & 1.24 $\pm$ 0.05 & 48.7 $\pm$ 2.7 \\
8 & 20.61 $\pm$ 0.29 & 2.14 $\pm$ 0.04 & 74.8 $\pm$ 2.8 \\
10 & 24.22 $\pm$ 0.31 & 3.44 $\pm$ 0.05 & 108.2 $\pm$ 3.1 \\
12 & 28.49 $\pm$ 0.28 & 4.94 $\pm$ 0.05 & 146.9 $\pm$ 2.9 \\
\bottomrule
\end{tabular}
\end{table}

\begin{figure}[t]
\centering
\includegraphics[width=0.98\linewidth]{fig4_scalability_density.pdf}
\caption{Decoupled fleet scalability evaluations: (a) Incremental $\text{D}^*$ Lite replanning latency and V2V mesh broadcast packets vs. dynamic crowd density ($N_{\text{humans}} \in [2..30]$, fixed fleet $N=4$); (b) Fleet makespan and single-file corridor mutex queueing wait time vs. autonomous fleet size ($N_{\text{carts}} \in [2..12]$, fixed crowd $N=10$).}
\label{fig:scalability}
\end{figure}

% ==============================================================================
% SECTION VII: CONCLUSION
% ==============================================================================
\section{Conclusion \& Future Research Directions}
This paper introduced the Distributed Dynamic Route Optimization ($\text{D}^2\text{RO}$) framework powered by Socially-Weighted Distributed Graph Optimization (SW-DGO). By formalizing a 5-component edge traversal cost function $C(u, v, t) = w_D D + w_M W_{\text{mesh}} + w_H H_{\text{prox}} + w_R R_{\text{lock}} + w_S S_{\text{trolley}}$, $\text{D}^2\text{RO}$ resolves the fundamental failure modes of prior Multi-Agent Path Finding and reactive navigation paradigms. Extensive Monte Carlo evaluations demonstrate a $100.0\%$ mission success rate, $0.00 \pm 0.00$ corridor deadlocks, $0.00 \pm 0.00$ intimate personal space violations, and sub-millisecond ($< 0.15\text{ ms}$) incremental $\text{D}^*$ Lite vertex repair times.

\subsection{Theoretical Guarantees \& Deadlock Prevention}
\begin{enumerate}
    \item \textbf{Heuristic Admissibility \& Incremental Optimality:} Because the Euclidean distance heuristic $h(s, s_{\text{goal}}) \le c^*(s, s_{\text{goal}})$ is strictly admissible and consistent, $\text{D}^*$ Lite guarantees optimal path extraction with respect to the currently observed edge cost field $C(u, v, t)$ while recomputing only the subgraph perturbed by dynamic obstacles \cite{koenig2002dstar}.
    \item \textbf{Deadlock Freedom in Single-File Bottlenecks:} In single-file passages where corridor width is insufficient for bidirectional passing, opposing agents entering simultaneously produce reactive live-locks ($\mathbf{F}_{\text{net}} = 0$). $\text{D}^2\text{RO}$ prevents deadlocks by assigning directional reservations and broadcasting `LOCK\_REQUEST` packets that dynamically inflate the reverse edge cost to $R_{\text{lock}} = \infty$. Opposing carts detect this infinity cost and immediately detour through parallel aisles or hold in Turnout Alcoves ($V_{\text{alcove}}$), empirically eliminating head-on deadlocks across all evaluated trials ($N_{\text{deadlock}} = 0.00 \pm 0.00$).
\end{enumerate}

\subsection{Real-World Physical Considerations \& Deployment Constraints}
\begin{enumerate}
    \item \textbf{RF Signal Attenuation \& Steel Fixture Multipath:} Metallic retail shelf gondolas and merchandise packaging attenuate $2.4\text{ GHz}$ RF signals by $-15\text{ to }-25\text{ dBm}$. $\text{D}^2\text{RO}$ addresses this through Sub-GHz ($868/915\text{ MHz}$) and dedicated short-range communications ($5.9\text{ GHz}$ IEEE 802.11p / DSRC) multi-hop TTL forwarding with autonomous exponential penalty decay ($W_{\text{mesh}}(t) = W_0 \cdot \exp(-\lambda t)$), ensuring that intermittent packet loss does not permanently poison the routing graph.
    \item \textbf{Indoor Positioning \& Sensor Fusion:} To mitigate wheel odometry drift on variable flooring (polished supermarket tiles, clinical vinyl, terminal carpets), ground truth localization is maintained via an Extended Kalman Filter (EKF) fusing $100\text{ Hz}$ wheel encoders, 6-DOF IMU gyros, Ultra-Wideband (UWB) transceiver trilateration (Decawave DWM1000), and 2D LiDAR scan-matching against architectural CAD floorplans.
    \item \textbf{Payload Invariance \& Non-Holonomic Kinodynamics:} Autonomous carts undergo substantial payload shifts ($15\text{ kg}$ empty to $65\text{ kg}$ fully loaded, a $+333\%$ mass increase), altering the center of gravity and rotational moment of inertia. $\text{D}^2\text{RO}$ modulates linear acceleration ($a_{\max}$) and steering curvature ($\kappa_{\max}$) based on real-time motor current draw and load-cell readings to prevent wheel scrubbing or tip-over during turns.
\end{enumerate}

\subsection{Future Research Directions}
\begin{enumerate}
    \item \textbf{Hybrid Learning-Guided Search:} Future investigations will explore coupling Graph Neural Networks (GNNs) or Deep Reinforcement Learning (e.g., PRIMAL / Learn to Follow) for global sub-goal allocation with the deterministic $\text{D}^2\text{RO}$ engine for micro-level kinodynamic path execution.
    \item \textbf{Heterogeneous Multi-Agent Ecosystems:} Expanding SW-DGO to mixed-fleet environments comprising autonomous delivery pods, heavy floor-scrubbing AGVs, and human-operated wheelchairs by introducing vehicle-specific agility weights into the $R_{\text{lock}}$ reservation tensor.
\end{enumerate}

\section*{Data and Code Availability}
To support scientific reproducibility and open science, all simulation source code, experimental benchmarking harnesses, verification test suites, raw CSV datasets ($N=100$ Monte Carlo trials across 5 algorithms, 5 cost configurations, and 3 cross-domain architectural environments), and high-resolution vector figures are available under the MIT open-source license at: \url{https://github.com/polla-fattah/D2RO}.

\section*{Acknowledgment}
The authors thank the Autonomous Systems and Robotics Laboratory for hardware access and computing support.

% ==============================================================================
% REFERENCES
% ==============================================================================
\bibliographystyle{IEEEtran}
\bibliography{references}

\end{document}
